\chapter{The grep command}

\section{What is grep?}

Grep stands for Global Regular Expression Print. We don't have to get into regex for the tutorial whatsoever.
You'll ussually see it used with pipes `$|$' but we haven't seen those yet so we'll stick with grep. What grep
does is that it looks for a specific string from a file or any output.

\lstinputlisting[
    language=bash,
    caption={10-bash.sh},
    captionpos=t,
    label={lst:10-bash},
    breaklines=true]
{../10-bash/10-bash.sh}

\section{Grep options}

First, we have the `-v' option which ussually means ``verbose'' or has to do with the command version.
In this case, it's like a not operator. It will output all instances where the specified word does not appear.

\begin{lstlisting}[
    language=bash, 
    caption={Terminal},
    captionpos=t,
    label={lst:10-terminal}, 
    mathescape=true, 
    breaklines=true]
grep -v Port /etc/ssh/ssh_config
\end{lstlisting}

The next one is very useful as it tells us at which line the specified word was found.

\begin{lstlisting}[
    language=bash, 
    caption={Terminal},
    captionpos=t,
    label={lst:10.1-terminal}, 
    mathescape=true, 
    breaklines=true]
grep -n Port /etc/ssh/ssh_config
\end{lstlisting}

Then we have the `-c' option for count, this will tell us how many instances of the specified word were
found

\begin{lstlisting}[
    language=bash, 
    caption={Terminal},
    captionpos=t,
    label={lst:10.2-terminal}, 
    mathescape=true, 
    breaklines=true]
grep -c Port /etc/ssh/ssh_config
\end{lstlisting}

If we wrote ``port'' instead of ``port'', it wouldn't have worked since grep is case sensitive. But we can
use the `-i' option or ignore case. very self explanatory

\begin{lstlisting}[
    language=bash, 
    caption={Terminal},
    captionpos=t,
    label={lst:10.3-terminal}, 
    mathescape=true, 
    breaklines=true]
grep -i port /etc/ssh/ssh_config
\end{lstlisting}

\pagebreak
\section{Search accross multiple files}

We can search accross more than a single file for any keyword.

\lstinputlisting[
    language=bash,
    caption={10-bash.sh},
    captionpos=t,
    label={lst:10-bash},
    breaklines=true]
{../10-bash/10.1-bash.sh}

The command will tell you the file name at which the instance of the keyword was found.\\ \\
There is also recursive search in case you want to search accross subdirectories:

\begin{lstlisting}[
    language=bash, 
    caption={Terminal},
    captionpos=t,
    label={lst:10.3-terminal}, 
    mathescape=true, 
    breaklines=true]
grep -r echo ./
\end{lstlisting}

\section{Combining options}

This is something general for all commands (I suppose). Instead of putting a billion flags in your command,
you can use a single flag and put all your options there.

\begin{lstlisting}[
    language=bash, 
    caption={Terminal},
    captionpos=t,
    label={lst:10.3-terminal}, 
    mathescape=true, 
    breaklines=true]
grep -nri error /var/log
\end{lstlisting}
